%************************************************
\chapter{Introduction}\label{ch:introduction}
%************************************************

Graph structures can be helpful as a tool in modeling several important domains.
Examples of applications of such models are found in domains as different as that of
predicting likely chemical bonds between atoms within medical science, and the problem
of fraud detection in the field of finance. The prevalence of graph-structured data
across many domains creates a need for efficient and accurate methods for modeling and
analyzing graphs.

\acsfont{GNN}s have emerged as a leading approach for representation learning on
graphs, combining the structural intuition behind classical graph algorithms with the
expressiveness and inductivity of neural networks. Unlike traditional graph
representation learning (\acsfont{GRL}) algorithms, \acsfont{GNN}s are inductive in
nature, and unlike most traditional neural networks they do not treat datapoints
independently of other inputs. This allows \acsfont{GNN}s to capture graph structure
as well as generalize to unseen nodes, making them highly effective when analyzing
graph-structured data. However, the graph structure also introduces a new challenge
absent from traditional neural networks: the challenge of growing training data. The
inter-connectivity between datapoints in a graph makes it impossible to consider
datapoints independently, as is typically assumed in stochastic gradient descent
(\acsfont{SGD}).

A natural solution is to store the graph in a dedicated graph database and query it
during training for the relevant neighbors of nodes in a mini-batch. Graph databases
are especially built for efficient graph traversal and feature look-up---precisely the
operations needed when training \acsfont{GNN}s on large graphs. Yet there have been
few studies exploring synergies between \acsfont{GNN}s and graph databases.

This thesis explores whether graph databases can be used as a scalable and efficient
solution for training \acsfont{GNN}s on large graphs. We build upon the
\texttt{PyTorch Geometric} machine learning library and couple it with subgraph
fetching from \texttt{Neo4j}, where the graph data resides.


%------------------------------------------------
\section{Research Questions}\label{sec:rq}
%------------------------------------------------

This thesis addresses the following research questions:

\begin{enumerate}
    \item \textbf{\acsfont{RQ}1 --- Bottlenecks:} What are the primary bottlenecks
    in \acsfont{GNN} training when using a graph database for graph storage, and to
    what extent can they be mitigated?

    \item \textbf{\acsfont{RQ}2 --- Sampling:} Which subgraph sampling strategies
    are most efficient when the graph is stored in a graph database?

    \item \textbf{\acsfont{RQ}3 --- Graph size:} For which graph sizes does a
    graph-database-based storage approach provide the greatest benefits?
\end{enumerate}


%------------------------------------------------
\section{Scope and Limitations}\label{sec:scope}
%------------------------------------------------

The experiments in this thesis are conducted on four datasets of varying sizes using a
two-layer Graph Convolutional Network (\acsfont{GCN}) for node classification. The
main emphasis is on the training pipeline, rather than on a particular downstream task
or model architecture. Since the core stages of the training pipeline remain broadly
similar across multiple graph learning tasks, the conclusions may have relevance beyond
the specific task considered here.
